% !TEX program = xelatex
\documentclass[12pt,a4paper]{ctexart}

% ============ 中文支持与字体 ============
% ctex 宏包会自动检测系统可用中文字体
% 若编译失败，可尝试显式指定：
%   \setCJKmainfont{SimSun}
%   \setCJKsansfont{SimHei}
%   \setCJKmonofont{FangSong}
% 或上传到 Overleaf 自动处理字体。

% ============ 数学与符号 ============
\usepackage{amsmath,amssymb,amsthm}
\usepackage{bm}
\usepackage{bbm}
\usepackage{mathtools}

% ============ 页面布局 ============
\usepackage[top=2.5cm, bottom=2.5cm, left=2.8cm, right=2.8cm]{geometry}
\usepackage{fancyhdr}
\pagestyle{fancy}
\fancyhf{}
\fancyhead[L]{第 2 周：Mirror Descent 与 AdaGrad 的 Regret 分析}
\fancyhead[R]{\thepage}
\fancyfoot[C]{上海财经大学·计算机机械社}
\renewcommand{\headrulewidth}{0.5pt}

% ============ 超链接与颜色 ============
\usepackage{xcolor}
\usepackage{hyperref}
\hypersetup{
  colorlinks=true,
  linkcolor=blue!60!black,
  citecolor=green!50!black,
  urlcolor=blue!60!black,
}
\definecolor{highlight}{rgb}{0.95, 0.95, 0.92}
\definecolor{codebg}{rgb}{0.97, 0.97, 0.97}

% ============ 代码环境 ============
\usepackage{listings}
\lstset{
  basicstyle=\ttfamily\small,
  backgroundcolor=\color{codebg},
  frame=single,
  framerule=0.5pt,
  rulecolor=\color{gray!40},
  breaklines=true,
  numbers=left,
  numbersep=6pt,
  numberstyle=\tiny\color{gray},
  tabsize=2,
  keywordstyle=\color{blue!70!black},
  commentstyle=\color{gray!50!black},
  stringstyle=\color{orange!70!black},
  showstringspaces=false,
}

% ============ 定理环境 ============
\theoremstyle{definition}
\newtheorem{definition}{定义}[section]
\newtheorem{example}{例}[section]
\newtheorem{exercise}{练习}[section]

\theoremstyle{plain}
\newtheorem{theorem}{定理}[section]
\newtheorem{lemma}[theorem]{引理}
\newtheorem{corollary}[theorem]{推论}
\newtheorem{proposition}[theorem]{命题}

\theoremstyle{remark}
\newtheorem{remark}{注记}[section]

% ============ 自定义命令 ============
\newcommand{\R}{\mathbb{R}}
\newcommand{\E}{\mathbb{E}}
\newcommand{\cL}{\mathcal{L}}
\newcommand{\cX}{\mathcal{X}}
\newcommand{\bbone}{\mathbbm{1}}
\newcommand{\argmin}{\operatornamewithlimits{argmin}}
\newcommand{\argmax}{\operatornamewithlimits{argmax}}
\newcommand{\diag}{\operatorname{diag}}
\newcommand{\Regret}{\mathrm{Regret}}
\newcommand{\prox}{\mathrm{prox}}

\DeclareMathOperator{\KL}{KL}

\title{\textbf{第 2 周：Mirror Descent 与 AdaGrad 的 Regret 分析}}
\author{上海财经大学·计算机机械社}
\date{2025–2026 学年（V2, \today）}

\begin{document}

\maketitle
\thispagestyle{fancy}

\begin{center}
\fbox{\parbox{0.85\textwidth}{\centering\small
\textbf{本周核心思想}\\[3pt]
「选优化器就是选几何。」\\
不同的 Bregman divergence 对应不同的几何结构，\\
Mirror Descent 通过在「镜像空间」中做梯度下降来利用这一几何，\\
而 AdaGrad 通过在线的对角预处理矩阵自动适应数据的几何。
}}
\end{center}

\vspace{6pt}

% ================================================================
\section{Bregman Divergence}
% ================================================================

\subsection{定义}

\begin{definition}[Bregman Divergence]
设 $\psi : \cX \to \R$ 是定义在闭凸集 $\cX \subseteq \R^d$ 上的一个严格凸且连续可微的函数。$\psi$ 诱导的 \textbf{Bregman 散度（Bregman Divergence）}定义为
\begin{equation}\label{eq:bregman-def}
  D_\psi(x, y) = \psi(x) - \psi(y) - \langle \nabla\psi(y),\, x - y \rangle.
\end{equation}
\end{definition}

\textbf{直观解释：} $D_\psi(x, y)$ 度量的是 $\psi(x)$ 与其在点 $y$ 处的一阶泰勒展开 $\psi(y) + \langle\nabla\psi(y), x-y\rangle$ 之间的差距。这正是在给定线性近似下，用 $\psi$ 来度量点 $x$ 与点 $y$ 之间「距离」的一种方式。

\begin{figure}[htbp]
\centering
\fbox{\parbox{0.8\textwidth}{\small
\textbf{Bregman 散度的几何解释}\\[4pt]
在点 $y$ 处对 $\psi$ 做一阶泰勒展开，得到支撑超平面。$D_\psi(x,y)$ 就是 $\psi(x)$ 与该超平面在 $x$ 处的垂直距离。由于 $\psi$ 是凸的，该距离始终非负。
}}
\end{figure}

\subsection{基本性质}

\begin{enumerate}
  \item \textbf{非负性：} $D_\psi(x, y) \ge 0$，且当 $\psi$ 严格凸时 $D_\psi(x, y) = 0 \iff x = y$。
  \item \textbf{非对称性：} 一般地 $D_\psi(x, y) \neq D_\psi(y, x)$——因此它是一种「散度」而非「距离」。
  \item \textbf{关于 $x$ 的凸性：} $\psi$ 的凸性保证了 $D_\psi(\cdot, y)$ 关于第一个变量是凸的。
  \item \textbf{线性变换下的等变性：} 若 $\widetilde{\psi}(x) = \psi(Ax)$，则 $D_{\widetilde{\psi}}(x, y) = D_\psi(Ax, Ay)$。
\end{enumerate}

\subsection{常见例子}

\begin{center}
\renewcommand{\arraystretch}{1.6}
\begin{tabular}{c|c|c}
  \hline
  \textbf{名称} & $\psi(x)$ & $D_\psi(x, y)$ \\
  \hline
  平方欧氏距离 & $\frac{1}{2}\|x\|_2^2$ & $\frac{1}{2}\|x-y\|_2^2$ \\
  \hline
  Mahalanobis 距离 & $\frac{1}{2}x^\top H x$（$H \succ 0$） & $\frac{1}{2}(x-y)^\top H (x-y)$ \\
  \hline
  KL 散度 & $\sum_i x_i \log x_i$（单纯形上） & $\sum_i x_i \log\frac{x_i}{y_i}$ \\
  \hline
  Itakura–Saito 距离 & $-\sum_i \log x_i$ & $\sum_i\!\big(\frac{x_i}{y_i} - \log\frac{x_i}{y_i} - 1\big)$ \\
  \hline
\end{tabular}
\end{center}

\textbf{关键洞察：} 不同的 $\psi$ 定义不同的「近邻结构」——即哪些点被认为是「近」的。这正是 Mirror Descent 可以适应不同几何结构的根本原因。

\subsection{三点恒等式 (Three-Point Identity)}

\begin{lemma}[三点恒等式]\label{lem:three-point}
设 $\psi$ 为严格凸且连续可微，则对任意 $x, y, z \in \cX$，有
\begin{equation}\label{eq:three-point}
  D_\psi(x, y) - D_\psi(x, z) + D_\psi(y, z)
  = \langle \nabla\psi(z) - \nabla\psi(y),\; x - y \rangle.
\end{equation}
\end{lemma}

\begin{proof}
展开左边：
\begin{align*}
  &\; \big[\psi(x) - \psi(y) - \langle\nabla\psi(y), x-y\rangle\big]
   - \big[\psi(x) - \psi(z) - \langle\nabla\psi(z), x-z\rangle\big] \\
  &\quad + \big[\psi(y) - \psi(z) - \langle\nabla\psi(z), y-z\rangle\big] \\[4pt]
  =&\; \cancel{\psi(x)} - \psi(y) - \langle\nabla\psi(y), x-y\rangle
   - \cancel{\psi(x)} + \psi(z) + \langle\nabla\psi(z), x-z\rangle \\
  &\quad + \psi(y) - \psi(z) - \langle\nabla\psi(z), y-z\rangle \\[4pt]
  =&\; \langle\nabla\psi(z), x - z\rangle - \langle\nabla\psi(y), x - y\rangle - \langle\nabla\psi(z), y - z\rangle \\[4pt]
  =&\; \langle\nabla\psi(z), x - y\rangle - \langle\nabla\psi(y), x - y\rangle \\[4pt]
  =&\; \langle \nabla\psi(z) - \nabla\psi(y),\; x - y \rangle.
\end{align*}
\end{proof}

\textbf{为什么重要？} 三点恒等式是 Mirror Descent regret 分析的核心代数工具。它将 Bregman 散度的变化与梯度差的线性内积联系起来，从而将「散度」语言翻译为「梯度」语言，使我们能够对 regret 进行 telescoping（望远镜求和）。

% ================================================================
\section{Mirror Descent 框架}
% ================================================================

\subsection{算法定义}

在传统的梯度下降中，我们在原始空间更新参数：
\[
  \theta_{t+1} = \theta_t - \eta g_t.
\]
这等价于求解近端问题：
\[
  \theta_{t+1}
  = \argmin_{\theta \in \cX} \Big\{
      \eta \langle g_t, \theta \rangle + \frac{1}{2}\|\theta - \theta_t\|_2^2
    \Big\}.
\]

\textbf{Mirror Descent (MD)} 将欧氏距离 $\frac{1}{2}\|\theta - \theta_t\|_2^2$ 替换为一般的 Bregman 散度 $D_\psi(\theta, \theta_t)$：

\begin{definition}[Mirror Descent 更新规则]
给定学习率 $\eta > 0$ 和 Bregman 势函数 $\psi$，
\begin{equation}\label{eq:md-update}
  \theta_{t+1}
  = \argmin_{\theta \in \cX} \Big\{
      \eta \langle g_t, \theta \rangle + D_\psi(\theta, \theta_t)
    \Big\}.
\end{equation}
\end{definition}

\subsection{从 MD 到预条件梯度下降}

\begin{theorem}[MD 退化为预条件 GD]\label{thm:md-to-gd}
取 $\psi(\theta) = \frac{1}{2}\theta^\top H \theta$，其中 $H \succ 0$（对称正定矩阵），
则 MD 更新退化为
\begin{equation}\label{eq:precond-gd}
  \theta_{t+1} = \theta_t - \eta H^{-1} g_t.
\end{equation}
\end{theorem}

\begin{proof}
代入 $\psi(\theta) = \frac{1}{2}\theta^\top H \theta$，有
\[
  D_\psi(\theta, \theta_t)
  = \frac{1}{2}\theta^\top H \theta
    - \frac{1}{2}\theta_t^\top H \theta_t
    - (H\theta_t)^\top (\theta - \theta_t)
  = \frac{1}{2}(\theta - \theta_t)^\top H (\theta - \theta_t).
\]

MD 子问题为
\[
  \min_{\theta \in \cX}
  \Big\{ \eta \langle g_t, \theta \rangle + \frac{1}{2}(\theta - \theta_t)^\top H (\theta - \theta_t) \Big\}.
\]

若 $\cX = \R^d$（无约束），则求导得最优条件：
\[
  \eta g_t + H(\theta - \theta_t) = 0
  \implies \theta = \theta_t - \eta H^{-1} g_t.
\]
\end{proof}

\textbf{解读：} $H$ 对梯度做了「预条件」（preconditioning）——在 $H$ 的特征值较大的方向（即 $\psi$ 在该方向更陡峭），步长被缩小；在 $H$ 的特征值较小的方向，步长被放大。通过明智地选择 $H$，我们可以做到适应数据几何的优化。

\subsection{镜像空间视角}

\begin{figure}[htbp]
\centering
\fbox{\parbox{0.85\textwidth}{\small
\textbf{「镜像」的直觉}\\[4pt]
MD 将参数 $\theta_t$ 通过 \textbf{镜像映射} $\nabla\psi$ 映射到对偶空间：
$y_t = \nabla\psi(\theta_t)$。在对偶空间中做标准的梯度下降：
$y_{t+1} = y_t - \eta g_t$。然后再通过 $(\nabla\psi)^{-1}$ 拉回原始空间：
$\theta_{t+1} = (\nabla\psi)^{-1}(y_{t+1})$。\\[4pt]
这就是「镜像」(Mirror) 的由来：原始空间的非欧更新，可以被理解为在对偶空间中的欧氏更新。
}}
\end{figure}

\textbf{关键的几何洞察：} $\psi$ 的选择决定了什么是一个「好的」局部距离度量。例如：
\begin{itemize}
  \item $\psi(x) = \frac{1}{2}\|x\|_2^2$ $\to$ 欧氏几何，适合稠密梯度场景；
  \item $\psi(x) = \sum_i x_i \log x_i$（负熵） $\to$ KL 几何，适合单纯形上的概率分布更新（如 EG 算法）；
  \item $\psi(\theta) = \frac{1}{2}\theta^\top H_t \theta$（$H_t$ 随时间变化） $\to$ AdaGrad 的思想雏形。
\end{itemize}

% ================================================================
\section{在线凸优化 (OCO) 与 Regret}
% ================================================================

\subsection{OCO 基本设定}

在在线凸优化中，学习器和环境进行 $T$ 轮的交互：

\begin{center}
\fbox{\parbox{0.75\textwidth}{\small
\textbf{每一轮 $t = 1, 2, \dots, T$：}\\[3pt]
1. 学习器选择一个决策 $\theta_t \in \cX$（凸集）。\\[2pt]
2. 环境（具有完全信息的对手）揭示一个凸损失函数 $\ell_t(\cdot)$。\\[2pt]
3. 学习器承受损失 $\ell_t(\theta_t)$ 并观察（子）梯度 $g_t \in \partial \ell_t(\theta_t)$。
}}
\end{center}

我们希望学习器的累计损失不比事后最优的固定决策差太多。

\begin{definition}[Regret（遗憾）]
学习器在 $T$ 轮后的 regret 定义为
\begin{equation}\label{eq:regret}
  \Regret_T = \sum_{t=1}^{T} \ell_t(\theta_t) - \min_{\theta \in \cX} \sum_{t=1}^{T} \ell_t(\theta).
\end{equation}
\end{definition}

\begin{itemize}
  \item $\Regret_T$ 是实际损失与事后最优固定决策的损失之差。
  \item 目标：设计算法使得 $\Regret_T = o(T)$，即 \textbf{平均 regret 趋近于零}。
  \item 更理想的目标：$\Regret_T = O(\sqrt{T})$。
\end{itemize}

\subsection{Mirror Descent 的 Regret Bound（通用版）}

\begin{theorem}[MD 通用 Regret Bound]\label{thm:md-regret}
设 $\psi$ 在 $\cX$ 上是 $\rho$-强凸的（关于范数 $\|\cdot\|$），学习率 $\eta > 0$。
则 Mirror Descent 的 regret 满足
\begin{equation}\label{eq:md-regret-bound}
  \Regret_T \le \frac{1}{\eta} D_\psi(\theta^*, \theta_1)
    + \frac{\eta}{2\rho} \sum_{t=1}^{T} \|g_t\|_*^2,
\end{equation}
其中 $\|\cdot\|_*$ 是 $\|\cdot\|$ 的对偶范数，$\theta^* = \argmin_{\theta\in\cX} \sum_t \ell_t(\theta)$。
\end{theorem}

\begin{proof}[证明（骨架）]
令 $\theta^*$ 为最优固定决策。由式~(\ref{eq:md-update}) 的最优条件：
\[
  \langle \eta g_t + \nabla\psi(\theta_{t+1}) - \nabla\psi(\theta_t),\; \theta^* - \theta_{t+1} \rangle \ge 0.
\]

利用三点恒等式（引理~\ref{lem:three-point}）消去 $\nabla\psi$，得到
\[
  \eta\langle g_t, \theta_t - \theta^*\rangle
  \le D_\psi(\theta^*, \theta_t) - D_\psi(\theta^*, \theta_{t+1})
    + \frac{\eta}{2\rho}\|g_t\|_*^2.
\]

对 $t=1$ 到 $T$ 求和，中间项 telescoping：
\[
  \sum_{t=1}^{T} \big[ D_\psi(\theta^*, \theta_t) - D_\psi(\theta^*, \theta_{t+1}) \big]
  = D_\psi(\theta^*, \theta_1) - D_\psi(\theta^*, \theta_{T+1})
  \le D_\psi(\theta^*, \theta_1).
\]

最后，由损失函数的凸性 $\ell_t(\theta_t) - \ell_t(\theta^*) \le \langle g_t, \theta_t - \theta^* \rangle$，即得结论。
\end{proof}

\textbf{该 bound 告诉我们什么？}
\begin{itemize}
  \item 第一项 $D_\psi(\theta^*, \theta_1)/\eta$ 度量初始化误差——与「最佳固定决策离初始点有多远」有关。
  \item 第二项 $(\eta/2\rho)\sum_t \|g_t\|_*^2$ 累积了梯度信息——反映了损失函数的「难度」。
  \item 最优学习率 $\eta = \Theta(1/\sqrt{T})$ 会给出 $\Regret_T = O(\sqrt{T})$。
\end{itemize}

% ================================================================
\section{AdaGrad}
% ================================================================

\subsection{动机：为什么需要自适应方法？}

回顾 SGD 的 regret bound（取 $\psi(x)=\frac{1}{2}\|x\|_2^2$，$\rho=1$，$\|\cdot\|_* = \|\cdot\|_2$）：
\[
  \Regret_T^{\mathrm{SGD}}
  \le \frac{D}{2\eta} + \frac{\eta}{2}\sum_{t=1}^{T} \|g_t\|_2^2.
\]
取最优 $\eta$ 后：
\[
  \Regret_T^{\mathrm{SGD}}
  = O\!\left(D \sqrt{\sum_{t=1}^{T} \|g_t\|_2^2}\right).
\]

\textbf{问题在哪？} 当梯度是稠密的（即每个坐标在每个时刻都有非零梯度），$\sum_t \|g_t\|_2^2 = \Theta(dT)$，于是
$\Regret_T^{\mathrm{SGD}} = O(\sqrt{dT})$。

但实际中梯度往往是稀疏的——例如只有 5\% 的坐标在每轮被激活。SGD 的 bound 并不能利用这种稀疏性，因为它对每个坐标施加了相同的有效学习率。$\|g_t\|_2^2$ 这个聚合指标掩盖了「哪些特征经常出现、哪些很少出现」的结构信息。

\textbf{AdaGrad 的核心思想：} 在每一轮，不仅要知道梯度的大小，还要知道 \textbf{每个坐标过去被更新的频率}，从而对活跃度高的特征使用较小的学习率，对活跃度低的特征使用较大的学习率。

\subsection{对角 AdaGrad 算法}

\begin{center}
\fbox{\parbox{0.80\textwidth}{
\textbf{算法 1：对角 AdaGrad (Diagonal AdaGrad)}\\[4pt]
\textbf{输入：} 初始参数 $\theta_1 \in \R^d$，常数 $\varepsilon > 0$。\\[2pt]
\textbf{初始化：} $s_0 = \bm{0} \in \R^d$（或 $s_{0,i} = 0, \forall i$）。\\[2pt]
\textbf{for} $t = 1, 2, \dots, T$ \textbf{do}\\[3pt]
\quad 1.\; 做预测 $\theta_t$。\\[1pt]
\quad 2.\; 接收损失函数 $\ell_t(\cdot)$ 并计算子梯度 $g_t \in \partial\ell_t(\theta_t)$。\\[1pt]
\quad 3.\; 更新累积平方梯度：
      $\displaystyle s_{t,i} = s_{t-1,i} + g_{t,i}^2,\quad \forall i=1,\dots,d$。\\[1pt]
\quad 4.\; 参数更新（每坐标自适应学习率）：
      $\displaystyle \theta_{t+1,i} = \theta_{t,i} - \frac{\eta}{\sqrt{s_{t,i} + \varepsilon}}\, g_{t,i},\quad \forall i$。\\[2pt]
\textbf{end for}
}}
\end{center}

\textbf{关键结构：}
\begin{itemize}
  \item 每个坐标 $i$ 拥有独立的学习率 $\eta_{t,i} = \eta / \sqrt{s_{t,i} + \varepsilon}$。
  \item $s_{t,i} = \sum_{k=1}^{t} g_{k,i}^2$ 是该坐标历史上所有梯度平方的累积和。
  \item 频繁更新的坐标（$s_{t,i}$ 大）$\to$ 学习率小；稀疏更新的坐标（$s_{t,i}$ 小）$\to$ 学习率大。
  \item $\varepsilon$ 用于避免除零。
\end{itemize}

\subsection{AdaGrad 的 Regret Bound 推导}

\begin{theorem}[对角 AdaGrad 的 Regret Bound]\label{thm:adagrad-bound}
设 $\ell_t$ 为凸函数，$\|g_t\|_\infty \le G_\infty$，$D_\infty = \max_{\theta\in\cX} \|\theta - \theta_1\|_\infty$。
则对角 AdaGrad 满足
\begin{equation}\label{eq:adagrad-main}
  \Regret_T \le \sqrt{2}\, D_\infty \sum_{i=1}^{d} \sqrt{\sum_{t=1}^{T} g_{t,i}^2}.
\end{equation}
\end{theorem}

\subsection*{推导步骤}

\textbf{第 1 步：MD 视角下的 AdaGrad。}

考虑时变的 Bregman 势函数：
\[
  \psi_t(\theta) = \frac{1}{2} \theta^\top H_t \theta,
  \quad \text{其中} \quad H_t = \frac{1}{\eta}\diag\big(\sqrt{s_{t,1}+\varepsilon},\; \dots,\; \sqrt{s_{t,d}+\varepsilon}\big).
\]

此时第 $t$ 步的 MD 更新为
\[
  \theta_{t+1} = \argmin_{\theta} \Big\{ \eta\langle g_t, \theta \rangle + D_{\psi_t}(\theta, \theta_t) \Big\}
            = \theta_t - \eta H_t^{-1} g_t,
\]
这正是 AdaGrad 的更新规则。

\textbf{第 2 步：单步 regret 分解。}

考虑损失函数为线性函数 $\ell_t(\theta) = \langle g_t, \theta \rangle$ 的情形（一般凸情形用一阶 Taylor 展开，结论一致）。我们有：
\begin{align}
  \langle g_t, \theta_t - \theta^*\rangle
  &\le D_{\psi_t}(\theta^*, \theta_t) - D_{\psi_t}(\theta^*, \theta_{t+1})
     + \frac{\eta}{2}\|g_t\|_{H_t^{-1}}^2 \label{eq:step-decomp} \\
  &= D_{\psi_t}(\theta^*, \theta_t) - D_{\psi_t}(\theta^*, \theta_{t+1})
     + \frac{\eta}{2} \sum_{i=1}^{d} \frac{g_{t,i}^2}{\sqrt{s_{t,i}+\varepsilon}}.
\end{align}

这使用了 $H_t^{-1} = \eta \cdot \diag\big(1/\sqrt{s_{t,1}+\varepsilon}, \dots, 1/\sqrt{s_{t,d}+\varepsilon}\big)$。

\textbf{第 3 步：Telescoping 的非平凡之处。}

由于 $\psi_t$ 随时间变化，$D_{\psi_t}(\theta^*, \theta_t)$ 和 $D_{\psi_{t-1}}(\theta^*, \theta_t)$ 不同，telescoping 不能简单地跨步进行。我们需要处理「漂移项」：
\begin{align*}
  &\sum_{t=1}^{T} \big[ D_{\psi_t}(\theta^*, \theta_t) - D_{\psi_t}(\theta^*, \theta_{t+1}) \big] \\
  =\;& D_{\psi_1}(\theta^*, \theta_1) - D_{\psi_T}(\theta^*, \theta_{T+1})
    + \sum_{t=1}^{T-1} \big[ D_{\psi_{t+1}}(\theta^*, \theta_{t+1}) - D_{\psi_t}(\theta^*, \theta_{t+1}) \big].
\end{align*}

对于 $\psi_t(\theta) = \frac{1}{2\eta}\theta^\top \diag(\sqrt{s_t + \varepsilon})\theta$，漂移项为
\[
  D_{\psi_{t+1}}(\theta^*, \theta_{t+1}) - D_{\psi_t}(\theta^*, \theta_{t+1})
  = \frac{1}{2\eta}\sum_{i=1}^{d}
    (\theta_i^* - \theta_{t+1,i})^2
    \big( \sqrt{s_{t+1,i}+\varepsilon} - \sqrt{s_{t,i}+\varepsilon} \big).
\]

\textbf{第 4 步：对 $T$ 步求和的最终 bound。}

综合第 2 步和第 3 步，经过细致的代数操作（详见 Duchi et al. 2011 附录 A），可以证明：
\[
  \sum_{t=1}^{T} \langle g_t, \theta_t - \theta^*\rangle
  \le \frac{1}{2\eta} \sum_{i=1}^{d} (\theta_i^* - \theta_{1,i})^2 \sqrt{s_{T,i} + \varepsilon}
     + \frac{\eta}{2} \sum_{t=1}^{T} \sum_{i=1}^{d} \frac{g_{t,i}^2}{\sqrt{s_{t,i}+\varepsilon}}.
\]

取 $\eta$ 最优，并使用不等式
\[
  \sum_{t=1}^{T} \frac{g_{t,i}^2}{\sqrt{s_{t,i} + \varepsilon}}
  \le 2\sqrt{s_{T,i} + \varepsilon},
\]
（这由 $\sum_t a_t / \sqrt{\sum_{k\le t} a_k} \le 2\sqrt{\sum_t a_t}$ 这一标准引理保证，其中 $a_t = g_{t,i}^2$），我们得到：
\[
  \Regret_T \le \sqrt{2}\, D_\infty \sum_{i=1}^{d} \sqrt{\sum_{t=1}^{T} g_{t,i}^2}.
\]

\subsection{为什么在稀疏梯度下远优于 SGD？}

将 AdaGrad 的 bound 与 SGD 的 bound 做对比。

\vspace{6pt}
\noindent
\fbox{\parbox{0.92\textwidth}{\small
\textbf{场景：} 假设每个样本只激活 $k \ll d$ 个坐标（例如 $k = 0.05d$），且非零梯度的值约为 $\pm 1$。
\begin{itemize}
  \item \textbf{SGD Bound:}
    $\displaystyle \Regret_T^{\mathrm{SGD}}
     = O\!\left(D_\infty \sqrt{\sum_{t=1}^T \sum_{i=1}^d g_{t,i}^2}\right)
     = O\!\left(D_\infty \sqrt{kT}\right).$
  \item \textbf{AdaGrad Bound:}
    $\displaystyle \Regret_T^{\mathrm{AdaGrad}}
     = O\!\left(D_\infty \sum_{i=1}^d \sqrt{\sum_{t=1}^T g_{t,i}^2}\right)
     = O\!\left(D_\infty \cdot d \cdot \sqrt{\frac{kT}{d}}\right)
     = O\!\left(D_\infty \sqrt{dkT}\right).$
\end{itemize}
等等——乍一看 AdaGrad 的 bound 似乎 \emph{更差}（因子 $\sqrt{dk}$ vs $\sqrt{k}$）。
}}

\textbf{关键在于 $D_\infty$ 的定义不同！} 上面的比较忽略了 $D_\infty$ 的量纲差异：
\begin{itemize}
  \item SGD 的 $D_\infty$ 是在 $\ell_2$ 意义下度量的，即 $D_\infty = \max \|\theta - \theta_1\|_2$。
  \item AdaGrad 的 $D_\infty$ 是在 $\ell_\infty$ 意义下度量的，即 $D_\infty = \max \|\theta - \theta_1\|_\infty$。
\end{itemize}

更精确的比较应当对 bound 取最优 trade-off。现在我们看最重要的洞察：

\begin{center}
\fbox{\parbox{0.88\textwidth}{\small
\textbf{稀疏梯度的优势：}\\[4pt]
考虑每个坐标 $i$ 在整个 $T$ 轮中被激活的总次数为 $\tau_i \ll T$。则
\[
  \sum_{i=1}^{d} \sqrt{\sum_{t=1}^{T} g_{t,i}^2}
  \approx \sum_{i=1}^{d} \sqrt{\tau_i}.
\]
而 SGD 的 bound 中，
\[
  \sqrt{\sum_{t=1}^{T} \|g_t\|_2^2}
  = \sqrt{\sum_{t=1}^{T} \sum_{i=1}^{d} g_{t,i}^2}
  \approx \sqrt{kT}.
\]

当 $\tau_i$ 极度不均匀时（$k \ll d$ 且某些特征主导），前者可以远小于后者。例如，若只有 $O(1)$ 个特征被频繁激活，其余特征的 $\tau_i = O(1)$，则
\[
  \sum_i \sqrt{\tau_i} = O(1 \cdot \sqrt{T} + d \cdot 1) = O(\sqrt{T} + d).
\]
相对地，SGD bound 总是 $O(\sqrt{dT})$——即使只有很少的特征是活跃的。

\textbf{因此，AdaGrad 在极端稀疏的数据（如文本分类中的 bag-of-words 特征）上能够实现远优于 SGD 的实际收敛。}
}}
\end{center}

\subsection{从「Regret Bound」到「实际收敛」的桥梁}
在 OCO 设定下，$\Regret_T/T \to 0$ 意味着平均损失趋近于最优固定决策的损失。对于随机优化（随机抽取 i.i.d.\ 数据），如果使用在线到批次的转换（online-to-batch conversion），AdaGrad 的 regret bound 直接给出：
\[
  \E[F(\bar{\theta}_T)] - F(\theta^*) = O\!\left(\frac{1}{T} \sum_{i=1}^{d} \E\!\left[\sqrt{\sum_{t=1}^{T}g_{t,i}^2}\right]\right),
\]
其中 $\bar{\theta}_T = \frac{1}{T}\sum_{t=1}^T \theta_t$。这使得 regret 分析成为算法性能保证的坚实基础。

% ================================================================
\section{FTRL (Follow-The-Regularized-Leader)}
% ================================================================

\subsection{基本概念}

Mirror Descent 不是唯一的在线学习框架。另一个核心框架是 \textbf{FTRL (Follow-The-Regularized-Leader)}。

\begin{definition}[FTRL 更新规则]
在第 $t$ 轮，FTRL 选择
\begin{equation}\label{eq:ftrl}
  \theta_{t+1}
  = \argmin_{\theta \in \cX}
    \left\{
      \sum_{k=1}^{t} \langle g_k, \theta \rangle + \frac{1}{\eta} R(\theta)
    \right\},
\end{equation}
其中 $R(\theta)$ 是一个正则化函数（通常是强凸的）。
\end{definition}

\textbf{直觉：} FTRL 直接选择那个在过去所有时刻的线性损失之和下表现最好的决策，同时用 $R(\theta)$ 做正则化以防止过拟合到有限的历史。

\subsection{MD 与 FTRL 的关系}

\begin{itemize}
  \item MD 是一种 \textbf{增量式} 方法——每步在上一点的基础上做局部更新。
  \item FTRL 是一种 \textbf{批量式} 方法——每步求解一个全局优化问题。
  \item 在特定条件下（线性损失 + 固定正则化器），MD 和 FTRL 是 \textbf{等价的}。更精确地说，若 $R(\theta) = \psi(\theta) - \psi(\theta_1)$（在合适的边界条件下），则 MD 产生的序列与 FTRL 产生的序列一致。
  \item 对于时变的正则化器（如 AdaGrad 中的 $H_t$），FTRL 的推广版 FTRL-Proximal 可以与 AdaGrad 取得类似的 regret bound。
\end{itemize}

\subsection{为什么引入 FTRL？}

\begin{enumerate}
  \item \textbf{统一视角：} AdaGrad 可以从「时变正则化器的 FTRL」角度推导出来。原始 AdaGrad 论文（Duchi et al.\ 2011）正是使用了 FTRL 框架。
  \item \textbf{稀疏性：} FTRL 配合 $\ell_1$ 正则化（尤其是 FTRL-Proximal）可以产生稀疏解，这对大规模在线学习（如 Google 的广告点击率预测）至关重要。
  \item \textbf{理论优雅性：} FTRL 的 regret 分析往往比 MD 更加直接——它不需要三点恒等式，而是直接使用凸函数的定义和最优条件。
\end{enumerate}

% ================================================================
\section{工程实践}
% ================================================================

\subsection{实现对角 AdaGrad}

以下 Python 实现展示了对角 AdaGrad 的核心逻辑：

\begin{lstlisting}[language=Python, caption={对角 AdaGrad 的 NumPy 实现}]
import numpy as np

class DiagonalAdaGrad:
    """对角 AdaGrad 优化器。

    参数
    ----------
    lr : float
        全局学习率 eta。
    eps : float
        平滑项，避免除零（默认 1e-8）。
    """
    def __init__(self, lr=0.1, eps=1e-8):
        self.lr = lr
        self.eps = eps
        self.s = None                # 累积平方梯度

    def update(self, theta, grad):
        """对参数执行一次更新的原地操作。

        参数
        ----------
        theta : ndarray, shape (d,)
            当前参数向量。
        grad : ndarray, shape (d,)
            当前批次的梯度。
        """
        if self.s is None:
            self.s = np.zeros_like(theta)

        # 累积平方梯度：s_{t,i} = s_{t-1,i} + g_{t,i}^2
        self.s += grad ** 2

        # 自适应更新：theta_{t+1,i} = theta_{t,i} - eta/sqrt(s_{t,i} + eps) * g_{t,i}
        theta -= self.lr / (np.sqrt(self.s) + self.eps) * grad
\end{lstlisting}

\subsection{构造稀疏梯度场景}

\begin{lstlisting}[language=Python, caption={稀疏梯度场景的构造}]
def generate_sparse_gradient_data(d, sparsity=0.05, n_samples=10000):
    """
    生成稀疏梯度场景的模拟数据。

    每个样本仅在随机选择的 5% 坐标上拥有非零梯度，
    以此模拟文本分类、推荐系统中常见的稀疏特征。

    参数
    ----------
    d : int
        数据维度。
    sparsity : float
        每个样本中非零特征的比例。
    n_samples : int
        样本数量。

    返回
    -------
    X : ndarray, shape (n_samples, d)
        特征矩阵。
    y : ndarray, shape (n_samples,)
        标签。
    theta_star : ndarray, shape (d,)
        真实的「最优」参数（用于生成标签）。
    """
    np.random.seed(42)
    theta_star = np.random.randn(d)       # 真实参数

    X = np.zeros((n_samples, d))
    for i in range(n_samples):
        # 随机选择 5% 的坐标作为活跃特征
        active = np.random.choice(d, size=int(d * sparsity), replace=False)
        X[i, active] = np.random.randn(len(active))

    # 生成带噪声的标签
    y = X @ theta_star + 0.1 * np.random.randn(n_samples)
    return X, y, theta_star
\end{lstlisting}

\subsection{比较 SGD vs AdaGrad}

\begin{lstlisting}[language=Python, caption={SGD vs AdaGrad 在稀疏场景的收敛比较}]
import matplotlib.pyplot as plt

def compare_optimizers_sparse(d=1000, sparsity=0.05,
                               n_samples=5000, epochs=20):
    """在稀疏梯度场景下比较 SGD 和 AdaGrad。"""
    X, y, theta_star = generate_sparse_gradient_data(
        d, sparsity, n_samples
    )

    # 初始化
    theta_sgd = np.zeros(d)
    theta_adagrad = np.zeros(d)
    adagrad = DiagonalAdaGrad(lr=0.1, eps=1e-8)

    loss_sgd, loss_adagrad = [], []

    for epoch in range(epochs):
        perm = np.random.permutation(n_samples)
        total_loss_sgd = 0.0
        total_loss_adagrad = 0.0

        for idx in perm:
            xi, yi = X[idx], y[idx]

            # --- SGD 更新 ---
            pred_sgd = np.dot(xi, theta_sgd)
            grad_sgd = 2 * (pred_sgd - yi) * xi
            theta_sgd -= 0.01 * grad_sgd    # 固定学习率
            total_loss_sgd += (pred_sgd - yi) ** 2

            # --- AdaGrad 更新 ---
            pred_ag = np.dot(xi, theta_adagrad)
            grad_ag = 2 * (pred_ag - yi) * xi
            adagrad.update(theta_adagrad, grad_ag)
            total_loss_adag += (pred_ag - yi) ** 2

        loss_sgd.append(np.sqrt(total_loss_sgd / n_samples))
        loss_adagrad.append(np.sqrt(total_loss_adag / n_samples))
        print(f"Epoch {epoch+1:2d}: "
              f"SGD RMSE={loss_sgd[-1]:.4f}, "
              f"AdaGrad RMSE={loss_adagrad[-1]:.4f}")

    # 绘图
    plt.figure(figsize=(8, 5))
    plt.plot(loss_sgd,  'o-', label='SGD (lr=0.01)',  alpha=0.8)
    plt.plot(loss_adagrad, 's-', label='AdaGrad (lr=0.1)', alpha=0.8)
    plt.xlabel('Epoch')
    plt.ylabel('RMSE')
    plt.title(f'稀疏梯度场景 (d={d}, sparsity={sparsity:.0%})')
    plt.legend()
    plt.grid(True, alpha=0.3)
    plt.tight_layout()
    plt.savefig('sgd_vs_adagrad_sparse.pdf', dpi=150)
    plt.show()

if __name__ == '__main__':
    compare_optimizers_sparse()
\end{lstlisting}

\textbf{预期结果：} 在稀疏梯度场景（$d = 1000$，仅 5\% 特征活跃）中，AdaGrad 的收敛速度应明显快于 SGD。SGD 对所有坐标使用相同的学习率，导致在稀疏特征上的有效步长过小；AdaGrad 为每个坐标自适应地调整学习率，使得稀疏坐标获得更大的步长。

\subsection{可视化每坐标的有效学习率}

\begin{lstlisting}[language=Python, caption={可视化每坐标有效学习率随时间的变化}]
def visualize_learning_rates(d=100, sparsity=0.05,
                              T=500):
    """
    可视化 AdaGrad 中每个坐标的有效学习率变化。

    在稀疏场景下，可以清楚地看到活跃坐标与非活跃坐标
    的学习率衰减速度的巨大差异。
    """
    np.random.seed(123)

    # 记录每个坐标的有效学习率
    lr_history = np.zeros((T, d))
    s = np.zeros(d)
    eta = 0.1
    eps = 1e-8

    for t in range(T):
        # 模拟稀疏梯度：每轮随机激活 5% 坐标
        grad = np.zeros(d)
        active = np.random.choice(
            d, size=int(d * sparsity), replace=False
        )
        grad[active] = np.random.randn(len(active))

        s += grad ** 2
        eff_lr = eta / (np.sqrt(s) + eps)
        lr_history[t] = eff_lr

    # 绘图
    fig, axes = plt.subplots(1, 2, figsize=(12, 5))

    # 左图：一些坐标的有效学习率轨迹
    ax = axes[0]
    # 选择活跃度不同的几个坐标
    coord_active = active[0]      # 持续活跃
    coord_rare = (active[0] + 50) % d  # 可能较不活跃
    coord_inactive = d - 1        # 几乎不活跃

    ax.semilogy(lr_history[:, coord_active],  label=f'坐标 {coord_active} (活跃)')
    ax.semilogy(lr_history[:, coord_rare],    label=f'坐标 {coord_rare}    (中等)')
    ax.semilogy(lr_history[:, coord_inactive], label=f'坐标 {coord_inactive} (不活跃)')
    ax.set_xlabel('轮数 t')
    ax.set_ylabel('有效学习率 (log scale)')
    ax.set_title('各坐标的有效学习率随时间衰减')
    ax.legend(fontsize=8)
    ax.grid(True, alpha=0.3)

    # 右图：最终时刻各坐标学习率的直方图
    ax = axes[1]
    ax.hist(lr_history[-1], bins=30, edgecolor='black',
            alpha=0.7)
    ax.set_xlabel('有效学习率')
    ax.set_ylabel('坐标数')
    ax.set_title(f'T={T} 时有效学习率的分布')
    ax.grid(True, alpha=0.3)

    plt.tight_layout()
    plt.savefig('adagrad_lr_visualization.pdf', dpi=150)
    plt.show()

if __name__ == '__main__':
    visualize_learning_rates()
\end{lstlisting}

\textbf{关键观察：}
\begin{itemize}
  \item 频繁更新的坐标的有效学习率迅速衰减（在 log 刻度下呈线性下降趋势）。
  \item 极少更新的坐标的学习率保持高位——只要它们出现，就能获得一次大幅的更新。
  \item 最终时刻的有效学习率分布呈长尾状——少量坐标学习率很小（经常被更新），大量坐标学习率接近初始值（很少被更新）。
\end{itemize}

% ================================================================
\section{检验标准}
% ================================================================

完成本周学习后，请确保能够：

\begin{enumerate}
  \item \textbf{Mirror Descent 几何：}
  阐明 Mirror Descent 如何通过改变 Bregman divergence 来改变优化几何。
  能够举出至少两个不同的 $\psi$ 及其对应的几何与适用场景。

  \item \textbf{退化推导：}
  手推当 $\psi(\theta) = \frac{1}{2}\theta^\top H \theta$ 时 MD 如何退化为预条件梯度下降。

  \item \textbf{AdaGrad Regret Bound 骨架：}
  能够写出对角 AdaGrad regret bound $R_T \le \sqrt{2} D_\infty \sum_{i=1}^d \sqrt{\sum_t g_{t,i}^2}$
  的推导骨架（不要求严格逐行写出所有代数，但要知道每一步用到了什么核心工具）：
  \begin{itemize}
    \item 单步 regret 分解（Bregman 散度 + 三点恒等式）
    \item 漂移项的处理（时变正则化器）
    \item 数学引理 $\sum_t a_t/\sqrt{\sum_{k\le t}a_k} \le 2\sqrt{\sum_t a_t}$
    \item 最优 $\eta$ 的选取与最终 bound 的合并
  \end{itemize}

  \item \textbf{稀疏优势：}
  能用自己的话解释为什么 AdaGrad 的 regret bound 在稀疏梯度场景下远优于 SGD 的 bound。

  \item \textbf{FTRL 概念：}
  能解释 FTRL 的基本思想，并说明其与 MD 的关系。
\end{enumerate}

% ================================================================
\section{思考题（选做）}
% ================================================================

\begin{exercise}[Bregman 散度的实验]
取 $\psi_1(x) = \frac{1}{2}x^2$ 和 $\psi_2(x) = x\log x$（在 $[0,1]$ 上）。
对于 $y=0.5$，绘制 $D_{\psi_1}(x, y)$ 与 $D_{\psi_2}(x, y)$ 随 $x$ 的变化曲线。
观察它们在不对称性、在边界附近的行为差异。
\end{exercise}

\begin{exercise}[AdaGrad Bound 的直接检验]
在 $d=100$ 的场景下，模拟 $T=1000$ 轮，随机生成稀疏梯度。计算并比较：
\[
  \sum_{t=1}^{T} \ell_t(\theta_t^{\mathrm{SGD}}) - \sum_{t=1}^{T} \ell_t(\theta^*),
  \qquad
  \sum_{t=1}^{T} \ell_t(\theta_t^{\mathrm{AdaGrad}}) - \sum_{t=1}^{T} \ell_t(\theta^*).
\]
验证实际 regret 是否逼近理论 bound 所预测的行为。
\end{exercise}

\begin{exercise}[FTRL-Proximal 的实现]
实现 FTRL-Proximal 算法并将其应用于与 AdaGrad 相同的稀疏场景。比较二者的稀疏性（解中零元素的比例）和收敛速度。
\end{exercise}

% ================================================================
\section*{参考文献}
% ================================================================

\begin{enumerate}
  \item \textbf{Duchi, J., Hazan, E., \& Singer, Y. (2011).}
  \textit{Adaptive Subgradient Methods for Online Learning and Stochastic Optimization.}
  Journal of Machine Learning Research, 12, 2121–2159.
  \begin{itemize}
    \item 这是 AdaGrad 的原始论文。定理 1–3 分别给出了全矩阵 AdaGrad、对角 AdaGrad 和复合目标下的 regret bound。
  \end{itemize}

  \item \textbf{Beck, A., \& Teboulle, M. (2003).}
  \textit{Mirror Descent and Nonlinear Projected Subgradient Methods for Convex Optimization.}
  Operations Research Letters, 31(3), 167–175.
  \begin{itemize}
    \item Mirror Descent 的经典论文，系统化了 Bregman divergence 在优化中的应用。
  \end{itemize}

  \item \textbf{McMahan, H. B. (2011).}
  \textit{Follow-the-Regularized-Leader and Mirror Descent: Equivalence Theorems and L1 Regularization.}
  In AISTATS.
  \begin{itemize}
    \item 系统证明了 MD 与 FTRL 之间的等价关系，并推广到 FTRL-Proximal。
  \end{itemize}

  \item \textbf{Hazan, E. (2016).}
  \textit{Introduction to Online Convex Optimization.}
  Foundations and Trends in Optimization, 2(3–4), 157–325.
  \begin{itemize}
    \item 在线凸优化的全面入门教材，涵盖 MD、FTRL 以及自适应方法。
  \end{itemize}
\end{enumerate}

\end{document}
